% Detailed security checklist with brands and menu paths. Compile with XeLaTeX.
\documentclass[10pt]{article}
\usepackage[letterpaper,margin=0.75in]{geometry}
\usepackage{fontspec}
\setmainfont{Libertinus Serif}
\setsansfont{Libertinus Sans}
\usepackage{xcolor,enumitem,amssymb,titlesec}
\definecolor{brownred}{HTML}{8B1A1A}
\titleformat{\section}{\large\sffamily\bfseries\color{brownred}}{}{0pt}{}
\titlespacing*{\section}{0pt}{10pt}{3pt}
\setlist[itemize]{label=$\square$,leftmargin=1.6em,itemsep=3pt,topsep=2pt,parsep=0pt}
\pagestyle{plain}
\setlength{\parindent}{0pt}
\raggedright
\begin{document}
{\LARGE\sffamily\bfseries\color{brownred} Detailed security checklist for professors}\par
\smallskip
\medskip
Start with level A. For a Brown-owned device or a Brown account, ask OIT first. We did not check the brands and the menu paths against the pages of the vendors or the reviewers. Check each menu path on your own device.

\section{A. This week (about 2 hours)}
\begin{itemize}
\item \textbf{Automatic updates. Restart today.}\\{\small iPhone: Settings $>$ General $>$ Software Update $>$ Automatic Updates.\\ Mac: System Settings $>$ General $>$ Software Update.\\ Windows: Settings $>$ Windows Update.\\ Chrome: menu $>$ Help $>$ About Google Chrome, then Relaunch.}
\item \textbf{Password manager: 1Password (paid), Bitwarden (free), or Dashlane.}\\{\small Built-in options: Apple Passwords, Google Password Manager.\\ Install it, write the main password on paper, and replace the passwords for email and the bank first.}
\item \textbf{Second sign-in step (2FA) on every account that offers one. Use a passkey where the site offers it.}\\{\small Google: go to g.co/passkeys $>$ Create a passkey.\\ Apple: turn on two-factor authentication and iCloud Keychain (Settings $>$ your name $>$ iCloud $>$ Passwords).\\ Other sites: look under Security or Sign-in in the account settings.\\ Strength, from most to least: security key, passkey, app code, text-message code.}
\item \textbf{Phone passcode of at least 6 digits.}\\{\small iPhone: Settings $>$ Face ID \& Passcode $>$ Stolen Device Protection.\\ Android: Settings $>$ Google $>$ All services $>$ Theft protection.}
\item \textbf{Carrier PIN or number lock, so that nobody can move your number to another phone (SIM swap).}\\{\small Verizon: Number Lock and SIM Protection.\\ T-Mobile: SIM Protection and Port Out Protection.\\ AT\&T: Wireless Account Lock (myAT\&T app $>$ Account).}
\item Habit: never approve a sign-in prompt, or enter a sign-in code, that you did not request.
\item Habit: never follow a web page's instruction to paste text into a terminal or the Windows Run box.
\item Habit: when someone asks for money or a password, call back on a number that you already have.
\item Family code word. Ask for it when a caller who sounds like a relative asks for money.
\end{itemize}
\section{B. This month (about 4 hours)}
\begin{itemize}
\item \textbf{Backup codes and a recovery contact, before you add security keys.}\\{\small Apple: Settings $>$ your name $>$ Sign-In \& Security $>$ Recovery Contacts.\\ Google: g.co/recovery-contacts, and print the backup codes.}
\item \textbf{Two security keys: Yubico Security Key C NFC (about \$29 each). Upgrade: YubiKey 5C NFC.}\\{\small Add both keys to each account, and keep the spare key at home.\\ Apple: Settings $>$ your name $>$ Sign-In \& Security $>$ Two-Factor Authentication $>$ Security Keys.\\ Google: Google Account $>$ Security $>$ 2-Step Verification.}
\item Remove text-message codes, once 2 other sign-in methods work.
\item \textbf{Browser.}\\{\small Chrome: Settings $>$ Privacy and security $>$ Security $>$ Always use secure connections.\\ Firefox: Settings $>$ Privacy \& Security $>$ HTTPS-Only Mode.\\ Content blocker: uBlock Origin (Firefox) or uBlock Origin Lite (Chrome, Edge, Safari).\\ Privacy extension: Privacy Badger.\\ Remove the extensions that you do not use.}
\item \textbf{Backup of your own files.}\\{\small Offline copy: an external drive with Time Machine (Mac) or File History (Windows).\\ Online backup: IDrive or Backblaze.\\ Restore 1 file as a test. Keep research and student data only on storage that Brown approves.}
\item \textbf{Disk encryption. Back up first, and save the recovery key away from the laptop.}\\{\small Mac: System Settings $>$ Privacy \& Security $>$ FileVault.\\ Windows: Settings $>$ Privacy \& security $>$ Device encryption (or BitLocker).}
\item \textbf{Credit freeze, free, at each bureau: Equifax, Experian, TransUnion.}\\{\small Lift it a few days before you apply for a loan or a card.}
\item Replace devices that no longer receive updates, such as an old home router.
\item \textbf{Delete the accounts that you no longer use.}\\{\small Find them in your password manager and in old ``welcome'' emails. Download the files that you need first.}
\end{itemize}
\section{C. When time allows, or at higher risk (personal accounts)}
\begin{itemize}
\item Google Advanced Protection Program: landing.google.com/advancedprotection.
\item \textbf{Apple Advanced Data Protection: Settings $>$ your name $>$ iCloud $>$ Advanced Data Protection.}\\{\small Set a recovery contact or a recovery key first.}
\item Apple Lockdown Mode: Settings $>$ Privacy \& Security $>$ Lockdown Mode.
\item Email masking: Hide My Email (Apple, with iCloud+), Firefox Relay, DuckDuckGo Email Protection, or Fastmail Masked Email.
\item Private Relay (Apple, with iCloud+): Settings $>$ your name $>$ iCloud $>$ Private Relay.
\item VPN: the Brown VPN for Brown resources. For privacy on other networks: Mullvad.
\item \textbf{Tracking.}\\{\small iPhone: Settings $>$ Privacy \& Security $>$ Tracking $>$ turn off Allow Apps to Request to Track.\\ Android: Settings $>$ Privacy $>$ Ads $>$ Delete advertising ID.\\ Data brokers: opt out yourself, or use Optery or EasyOptOuts (about \$20 per year).}
\end{itemize}
\section{Research code}
\begin{itemize}
\item \textbf{Wait about 7 days before you install a new release of an open-source code package (R, Python, npm).}\\{\small Update phones, laptops, and browsers at once.}
\item Run AI coding agents in a container.
\item API keys: keep them out of code and out of prompts. Give each project its own key, with the fewest permissions and a spending limit.
\item For a key on a fixed server, limit the key to the address of that server, where the provider offers this setting. Delete unused keys.
\end{itemize}
\section{After a suspected break-in}
\begin{itemize}
\item Call OIT first, from another phone. Follow its instructions for the device.
\item Do not erase the device, delete files, or try to repair it.
\item On a different device, change the password of the affected account, and sign out of all other devices.
\item If money moved, call the bank.
\end{itemize}

\section{Sources of brand names}
{\small 1Password and Bitwarden: Wirecutter's picks, according to third-party show notes. Dashlane: the speaker's recommendation. uBlock Origin and Privacy Badger: Wirecutter's recommendations, according to the speaker. Yubico keys: Wirecutter's pick, according to a Yubico-hosted copy of the Wirecutter article. IDrive: Wirecutter's pick, according to IDrive's own blog (January 2026). Backblaze: the pick of the review site Cloudwards. Mullvad: the top pick of the review site Rtings. Optery and EasyOptOuts: the 2 services that removed the most profiles in a Consumer Reports test of August 2024.}
\end{document}
